\newpage

\section{线性代数综合复习}

本章将线性代数七章内容融会贯通，通过跨章节综合题、提高题和证明题，帮助读者建立完整的知识体系。主体题目服务于考研数学一主线综合训练；少量使用 Cayley--Hamilton 定理、广义特征值或 Sylvester 不等式的题目标为“拓展”，不纳入主线备考要求。

\subsection{基础综合题}

\begin{enumerate}

\item \textbf{（行列式+矩阵）} 设 \(A = \begin{pmatrix} 1 & 2 & 1 \\ 2 & 5 & 2 \\ 1 & 2 & 3 \end{pmatrix}\)，\(B\) 为 \(3\) 阶方阵且 \(|B| = 2\)，计算 \(|2A^*B^{-1}A^{\mathsf{T}}|\)。

\ifshowsolutions%
\boxed{\textbf{解}}
先求 \(|A|\)：
\[
|A| = \begin{vmatrix}
1 & 2 & 1 \\
2 & 5 & 2 \\
1 & 2 & 3
\end{vmatrix}
\xrightarrow{r_2-2r_1,\ r_3-r_1}
\begin{vmatrix}
1 & 2 & 1 \\
0 & 1 & 0 \\
0 & 0 & 2
\end{vmatrix} = 1 \cdot 1 \cdot 2 = 2.
\]

利用行列式性质：
\[
|A^*| = |A|^{n-1} = 4,\quad
|A^{\mathsf{T}}| = |A| = 2,\quad
|B^{-1}| = |B|^{-1} = \frac{1}{2}.
\]

故
\[
|2A^*B^{-1}A^{\mathsf{T}}|
= 2^3 \cdot |A^*| \cdot |B^{-1}| \cdot |A^{\mathsf{T}}|
= 8 \cdot 4 \cdot \frac{1}{2} \cdot 2 = 32.
\]

\fi%
\item \textbf{（行列式+矩阵）} 设 \(n\) 阶矩阵 \(A\) 满足 \(A^2 - 3A + 2I = O\)。求 \(|A - I|\) 与 \(|A + I|\) 的所有可能取值。

\ifshowsolutions%
\boxed{\textbf{解}}
由 \(A^2 - 3A + 2I = O\) 得
\[
(A - I)(A - 2I) = O.
\]

取行列式：\(|A - I| \cdot |A - 2I| = 0\)。

设 \(\lambda\) 为 \(A\) 的任一特征值，则由 \((A - I)(A - 2I) = O\) 代入得 \((\lambda - 1)(\lambda - 2) = 0\)，故 \(\lambda \in \{1, 2\}\)。
\(A - I\) 的特征值为 \(\lambda - 1 \in \{0, 1\}\)。因此
\begin{itemize}
    \item 若 \(A\) 有特征值 \(1\)，则 \(|A - I| = 0\)；
    \item 若 \(A\) 没有特征值 \(1\)，则 \(A-I\) 可逆。由 \((A-I)(A-2I)=O\) 两边左乘 \((A-I)^{-1}\)，得 \(A=2I\)，此时 \(|A-I|=|I|=1\)。
\end{itemize}

\(A + I\) 的特征值为 \(\lambda + 1 \in \{2, 3\}\)。设 \(A\) 的特征值中恰有 \(k\) 个为 \(1\)，\(n-k\) 个为 \(2\)，则
\[
|A + I| = 2^k \cdot 3^{n-k},\quad k = 0, 1, \ldots, n.
\]
故 \(|A + I|\) 的所有可能取值为 \(\{2^k \cdot 3^{n-k} \mid k = 0, 1, \ldots, n\}\)。

\boxed{\text{注}} 此题将矩阵方程与特征值理论结合，核心在于将矩阵方程转化为特征值的约束条件。

\fi%
\item \textbf{（向量+方程组）} 设 \(\boldsymbol{\alpha}_1 = (1, 2, -1, 0)^{\mathsf{T}}\)，\(\boldsymbol{\alpha}_2 = (2, 5, a, 1)^{\mathsf{T}}\)，\(\boldsymbol{\alpha}_3 = (1, 1, -4, b)^{\mathsf{T}}\)，\(\boldsymbol{\beta} = (3, 8, 0, c)^{\mathsf{T}}\)。问 \(a, b, c\) 取何值时，

\begin{enumerate}[label=(\roman*)]
    \item \(\boldsymbol{\beta}\) 可由 \(\boldsymbol{\alpha}_1, \boldsymbol{\alpha}_2, \boldsymbol{\alpha}_3\) 唯一线性表示？
    \item \(\boldsymbol{\beta}\) 可由 \(\boldsymbol{\alpha}_1, \boldsymbol{\alpha}_2, \boldsymbol{\alpha}_3\) 线性表示但表示法不唯一？
    \item \(\boldsymbol{\beta}\) 不能由 \(\boldsymbol{\alpha}_1, \boldsymbol{\alpha}_2, \boldsymbol{\alpha}_3\) 线性表示？
\end{enumerate}

\ifshowsolutions%
\boxed{\textbf{解}}
该问题等价于非齐次线性方程组
\[
x_1\boldsymbol{\alpha}_1 + x_2\boldsymbol{\alpha}_2 + x_3\boldsymbol{\alpha}_3 = \boldsymbol{\beta}
\]
的解的存在性与唯一性问题。写出增广矩阵并作初等行变换：
\[
(A \mid \boldsymbol{\beta}) =
\begin{pmatrix}
1 & 2 & 1 & 3 \\
2 & 5 & 1 & 8 \\
-1 & a & -4 & 0 \\
0 & 1 & b & c
\end{pmatrix}
\xrightarrow{r_2-2r_1,\ r_3+r_1}
\begin{pmatrix}
1 & 2 & 1 & 3 \\
0 & 1 & -1 & 2 \\
0 & a+2 & -3 & 3 \\
0 & 1 & b & c
\end{pmatrix}
\]
\[
\xrightarrow{r_3-(a+2)r_2,\ r_4-r_2}
\begin{pmatrix}
1 & 2 & 1 & 3 \\
0 & 1 & -1 & 2 \\
0 & 0 & a-1 & -2a-1 \\
0 & 0 & b+1 & c-2
\end{pmatrix}.
\]

分类讨论：
\begin{enumerate}[label=(\roman*)]
    \item 当 \(a \neq 1\) 时，前三行系数矩阵满秩（秩 \(=3\)）。此时还需第四行与前三行相容：将第三行消去第四行的 \(b+1\) 位置，得
    \[
    \begin{pmatrix}
    1 & 2 & 1 & 3 \\
    0 & 1 & -1 & 2 \\
    0 & 0 & a-1 & -2a-1 \\
    0 & 0 & 0 & c-2 - \frac{(b+1)(-2a-1)}{a-1}
    \end{pmatrix}.
    \]
    故当且仅当 \(c-2 = \dfrac{(b+1)(-2a-1)}{a-1}\) 时方程组有唯一解；否则无解。
    \item 当 \(a = 1\) 时，第三行化为 \((0, 0, 0, -3)\)，矛盾！故此时无论 \(b,c\) 为何值，均无解。
\fi%
\end{enumerate}

综上：(i) \(a \neq 1\) 且满足相容条件时唯一表示；(ii) 不可能出现无穷多解，因为所有相容情形均有 \(a\neq1\)，此时系数矩阵秩为 \(3\)；(iii) 其余情况均不可表示。需要注意，\(a=1\) 时方程组恒不相容，而向量组本身是否相关还取决于 \(b\)：仅当 \(b=-1\) 时系数矩阵秩降为 \(2\)。

\item \textbf{（向量+方程组）} 已知 \(A\) 为 \(m \times n\) 实矩阵，\(\boldsymbol{0}\ne\boldsymbol{b} \in \mathbb{R}^m\)。若 \(A\boldsymbol{x} = \boldsymbol{0}\) 的基础解系为 \(\boldsymbol{\xi}_1, \boldsymbol{\xi}_2\)，\(\boldsymbol{\eta}_0\) 为 \(A\boldsymbol{x} = \boldsymbol{b}\) 的一个特解。证明：

\begin{enumerate}[label=(\roman*)]
    \item \(\boldsymbol{\eta}_0,\ \boldsymbol{\eta}_0 + \boldsymbol{\xi}_1,\ \boldsymbol{\eta}_0 + \boldsymbol{\xi}_2\) 线性无关；
    \item \(A\boldsymbol{x} = \boldsymbol{b}\) 的任意 \(4\) 个解必线性相关。
\end{enumerate}

\ifshowsolutions%
\boxed{\textbf{解}}
(i) 设
\[
k_0\boldsymbol{\eta}_0 + k_1(\boldsymbol{\eta}_0 + \boldsymbol{\xi}_1) + k_2(\boldsymbol{\eta}_0 + \boldsymbol{\xi}_2) = \boldsymbol{0},
\]
整理得
\[
(k_0+k_1+k_2)\boldsymbol{\eta}_0 + k_1\boldsymbol{\xi}_1 + k_2\boldsymbol{\xi}_2 = \boldsymbol{0}. \tag{1}
\]

在 (1) 两端左乘 \(A\)，利用 \(A\boldsymbol{\eta}_0 = \boldsymbol{b}\)，\(A\boldsymbol{\xi}_1 = A\boldsymbol{\xi}_2 = \boldsymbol{0}\) 得
\[
(k_0+k_1+k_2)\boldsymbol{b} = \boldsymbol{0}.
\]
因为题设给出 \(\boldsymbol{b} \neq \boldsymbol{0}\)，故
\[
k_0 + k_1 + k_2 = 0.
\]

代回 (1)：\(k_1\boldsymbol{\xi}_1 + k_2\boldsymbol{\xi}_2 = \boldsymbol{0}\)。由基础解系中向量的线性无关性知 \(k_1 = k_2 = 0\)，进而 \(k_0 = 0\)。故三向量线性无关。

(ii) \(A\boldsymbol{x} = \boldsymbol{b}\) 的通解为
\[
\boldsymbol{\eta}_0+c_1\boldsymbol{\xi}_1+c_2\boldsymbol{\xi}_2.
\]
因此每个解都属于线性子空间
\(\operatorname{span}\{\boldsymbol{\eta}_0,\boldsymbol{\xi}_1,\boldsymbol{\xi}_2\}\)，其维数至多为 \(3\)，故任意 \(4\) 个解必线性相关。

\boxed{\text{注}} “解集是二维仿射子空间”不能推出任意三个解线性相关；当该仿射平面不经过原点时，其中完全可能存在三个线性无关向量。(i) 正是一个例子。

\fi%
\item \textbf{（特征值+二次型）} 设三阶实对称矩阵 \(A\) 的特征值为 \(\lambda_1 = 2\)，\(\lambda_2 = \lambda_3 = -1\)，且对应于 \(\lambda_1 = 2\) 的特征向量为 \(\boldsymbol{p}_1 = (1, 1, 1)^{\mathsf{T}}\)。

\begin{enumerate}[label=(\roman*)]
    \item 求矩阵 \(A\)；
    \item 写出二次型 \(f(\boldsymbol{x}) = \boldsymbol{x}^{\mathsf{T}} A \boldsymbol{x}\) 在正交变换下的标准形。
\end{enumerate}

\ifshowsolutions%
\boxed{\textbf{解}}
(i) 实对称矩阵属于不同特征值的特征向量正交，因此 \(\lambda_2 = \lambda_3 = -1\) 对应的特征向量均与 \(\boldsymbol{p}_1 = (1,1,1)^{\mathsf{T}}\) 正交。求解 \(x_1+x_2+x_3 = 0\)，取基础解系：
\[
\boldsymbol{p}_2 = (1, -1, 0)^{\mathsf{T}},\quad \boldsymbol{p}_3 = (1, 0, -1)^{\mathsf{T}}.
\]

对 \(\boldsymbol{p}_2, \boldsymbol{p}_3\) 进行 Schmidt 正交化：
\[
\boldsymbol{q}_2 = \boldsymbol{p}_2 = (1, -1, 0)^{\mathsf{T}},
\]
\[
\boldsymbol{q}_3 = \boldsymbol{p}_3 - \frac{\boldsymbol{p}_3 \cdot \boldsymbol{q}_2}{\|\boldsymbol{q}_2\|^2}\boldsymbol{q}_2
= (1, 0, -1)^{\mathsf{T}} - \frac{1}{2}(1, -1, 0)^{\mathsf{T}}
= \left(\frac{1}{2}, \frac{1}{2}, -1\right)^{\mathsf{T}}.
\]

单位化：
\[
\boldsymbol{e}_1 = \frac{1}{\sqrt{3}}(1, 1, 1)^{\mathsf{T}},\quad
\boldsymbol{e}_2 = \frac{1}{\sqrt{2}}(1, -1, 0)^{\mathsf{T}},\quad
\boldsymbol{e}_3 = \frac{1}{\sqrt{6}}(1, 1, -2)^{\mathsf{T}}.
\]

正交矩阵 \(Q = (\boldsymbol{e}_1, \boldsymbol{e}_2, \boldsymbol{e}_3)\)，则
\[
A = Q \begin{pmatrix} 2 & 0 & 0 \\ 0 & -1 & 0 \\ 0 & 0 & -1 \end{pmatrix} Q^{\mathsf{T}}.
\]

计算：
\begin{align*}
A &= 2\boldsymbol{e}_1\boldsymbol{e}_1^{\mathsf{T}} - \boldsymbol{e}_2\boldsymbol{e}_2^{\mathsf{T}} - \boldsymbol{e}_3\boldsymbol{e}_3^{\mathsf{T}} \\[4pt]
  &= \frac{2}{3}\begin{pmatrix}1\\1\\1\end{pmatrix}(1,1,1)
     - \frac{1}{2}\begin{pmatrix}1\\-1\\0\end{pmatrix}(1,-1,0)
     - \frac{1}{6}\begin{pmatrix}1\\1\\-2\end{pmatrix}(1,1,-2) \\[4pt]
  &= \begin{pmatrix}
     0 & 1 & 1 \\ 1 & 0 & 1 \\ 1 & 1 & 0
     \end{pmatrix}.
\end{align*}

(ii) 在正交变换 \(\boldsymbol{x} = Q\boldsymbol{y}\) 下：
\[
f = \boldsymbol{x}^{\mathsf{T}} A \boldsymbol{x}
  = \boldsymbol{y}^{\mathsf{T}} (Q^{\mathsf{T}} A Q) \boldsymbol{y}
  = 2y_1^2 - y_2^2 - y_3^2.
\]

\fi%
\item \textbf{（特征值+二次型）} 判断二次型 \(f(x_1, x_2, x_3) = 2x_1^2 + 2x_2^2 + 3x_3^2 + 2x_1x_2 + 2x_1x_3 + 2x_2x_3\) 的正定性，并写出其规范形。

\ifshowsolutions%
\boxed{\textbf{解}}
二次型矩阵为
\[
A = \begin{pmatrix}
2 & 1 & 1 \\
1 & 2 & 1 \\
1 & 1 & 3
\end{pmatrix}.
\]

计算顺序主子式：
\[
\Delta_1 = 2 > 0,\quad
\Delta_2 = \begin{vmatrix} 2 & 1 \\ 1 & 2 \end{vmatrix} = 3 > 0,
\]
\[
\Delta_3 = |A| = 2\begin{vmatrix} 2 & 1 \\ 1 & 3 \end{vmatrix}
- 1\begin{vmatrix} 1 & 1 \\ 1 & 3 \end{vmatrix}
+ 1\begin{vmatrix} 1 & 2 \\ 1 & 1 \end{vmatrix}
= 2 \cdot 5 - 2 + (-1) = 7 > 0.
\]

所有顺序主子式均大于零，故 \(A\) 正定，即二次型正定。正定二次型的规范形为
\[
f = z_1^2 + z_2^2 + z_3^2.
\]

\boxed{\text{注}} 另法：对 \(A\) 作合同变换（配方法）：
\begin{align*}
f &= 2x_1^2 + 2x_2^2 + 3x_3^2 + 2x_1x_2 + 2x_1x_3 + 2x_2x_3 \\
  &= 2\left(x_1 + \frac{1}{2}x_2 + \frac{1}{2}x_3\right)^2
     + \frac{3}{2}\left(x_2 + \frac{1}{3}x_3\right)^2
     + \frac{7}{3}x_3^2 > 0 \ (\boldsymbol{x} \neq \boldsymbol{0}).
\end{align*}
同样证得正定性。

\fi%
\end{enumerate}

\subsection{提高综合题}

\begin{enumerate}

\item \textbf{（分块矩阵+行列式）} 设 \(A, B\) 均为 \(n\) 阶方阵，证明：
\[
\begin{vmatrix} A & B \\ B & A \end{vmatrix} = |A + B| \cdot |A - B|.
\]
当 \(A = \begin{pmatrix} 3 & 1 \\ 2 & 4 \end{pmatrix}\)，\(B = \begin{pmatrix} 1 & 2 \\ 3 & 1 \end{pmatrix}\) 时，计算分块行列式的值。

\ifshowsolutions%
\boxed{\textbf{解}}
利用分块矩阵的初等变换（打洞法）：
\[
\begin{aligned}
&\begin{pmatrix} I & I \\ O & I \end{pmatrix}
\begin{pmatrix} A & B \\ B & A \end{pmatrix}
=\begin{pmatrix} A+B & A+B \\ B & A \end{pmatrix},\\
&\begin{pmatrix} A+B & A+B \\ B & A \end{pmatrix}
\begin{pmatrix} I & -I \\ O & I \end{pmatrix}
=\begin{pmatrix} A+B & O \\ B & A-B \end{pmatrix}.
\end{aligned}
\]

两边取行列式。初等块变换矩阵 \(\begin{pmatrix} I & I \\ O & I \end{pmatrix}\) 和 \(\begin{pmatrix} I & -I \\ O & I \end{pmatrix}\) 的行列式均为 \(1\)（分块上三角矩阵，对角块均为 \(I\)），故
\[
\begin{vmatrix} A & B \\ B & A \end{vmatrix}
= \begin{vmatrix} A+B & O \\ B & A-B \end{vmatrix}
= |A+B| \cdot |A-B|.
\]

计算具体数值：
\[
A+B = \begin{pmatrix} 4 & 3 \\ 5 & 5 \end{pmatrix},\ |A+B| = 20 - 15 = 5,
\]
\[
A-B = \begin{pmatrix} 2 & -1 \\ -1 & 3 \end{pmatrix},\ |A-B| = 6 - 1 = 5.
\]
故
\[
\begin{vmatrix} A & B \\ B & A \end{vmatrix} = 5 \times 5 = 25.
\]

\boxed{\text{注}} 「打洞法」是处理分块矩阵行列式的高频技巧，其本质是通过可逆行（列）变换将分块矩阵化为分块上（下）三角形式。

\fi%
\item \textbf{（拓展：秩方程与 Sylvester 不等式）} 设 \(A\) 为 \(m \times n\) 矩阵，\(B\) 为 \(n \times s\) 矩阵。证明：
\[
\operatorname{rank}(AB) = \operatorname{rank}(B) - \dim(\ker A \cap \operatorname{col}(B)).
\]
并由此推出 \(\operatorname{rank}(AB) \geq \operatorname{rank}(A) + \operatorname{rank}(B) - n\)（Sylvester 不等式）。

\ifshowsolutions%
\boxed{\textbf{解}}
考虑线性映射 \(B: \mathbb{R}^s \to \mathbb{R}^n\)，\(A: \mathbb{R}^n \to \mathbb{R}^m\)。考察 \(A\) 在 \(\operatorname{col}(B)\) 上的限制 \(A|_{\operatorname{col}(B)}: \operatorname{col}(B) \to \mathbb{R}^m\)。

由秩—零化度定理：
\[
\dim(\operatorname{col}(B)) = \dim(A(\operatorname{col}(B))) + \dim(\ker A \cap \operatorname{col}(B)).
\]

而 \(\dim(\operatorname{col}(B)) = \operatorname{rank}(B)\)，\(A(\operatorname{col}(B)) = \operatorname{col}(AB)\)，故 \(\dim(A(\operatorname{col}(B))) = \operatorname{rank}(AB)\)。因此
\[
\operatorname{rank}(B) = \operatorname{rank}(AB) + \dim(\ker A \cap \operatorname{col}(B)),
\]
即
\[
\operatorname{rank}(AB) = \operatorname{rank}(B) - \dim(\ker A \cap \operatorname{col}(B)).
\]

对于 Sylvester 不等式：由 \(\ker A \cap \operatorname{col}(B) \subseteq \ker A\) 知
\[
\dim(\ker A \cap \operatorname{col}(B)) \leq \dim(\ker A) = n - \operatorname{rank}(A).
\]
代入上式：
\[
\operatorname{rank}(AB) \geq \operatorname{rank}(B) - (n - \operatorname{rank}(A)) = \operatorname{rank}(A) + \operatorname{rank}(B) - n.
\]

等号成立当且仅当 \(\ker A \subseteq \operatorname{col}(B)\)。

\boxed{\text{注}} Sylvester 不等式是线性代数中最重要的秩不等式之一，与 Frobenius 不等式 \(\operatorname{rank}(AB) + \operatorname{rank}(BC) \leq \operatorname{rank}(B) + \operatorname{rank}(ABC)\) 并列。

\fi%
\item \textbf{（含参方程组+特征值）} 设 \(A = \begin{pmatrix} 1 & a & 1 \\ a & 1 & -1 \\ 1 & -1 & 1 \end{pmatrix}\)，\(\boldsymbol{b} = \begin{pmatrix} 1 \\ 1 \\ -1 \end{pmatrix}\)。

\begin{enumerate}[label=(\roman*)]
    \item 讨论 \(a\) 取何值时，方程组 \(A\boldsymbol{x} = \boldsymbol{b}\) 有唯一解、无穷多解、无解；
    \item 求 \(A\) 的特征值与 \(a\) 的关系，并讨论 \(A\) 何时可对角化。
\end{enumerate}

\ifshowsolutions%
\boxed{\textbf{解}}
(i) 写出增广矩阵并作行变换：
\[
(A \mid \boldsymbol{b}) = \begin{pmatrix}
1 & a & 1 & 1 \\
a & 1 & -1 & 1 \\
1 & -1 & 1 & -1
\end{pmatrix}
\xrightarrow{r_3 \leftrightarrow r_1}
\begin{pmatrix}
1 & -1 & 1 & -1 \\
a & 1 & -1 & 1 \\
1 & a & 1 & 1
\end{pmatrix}
\]
\[
\xrightarrow{r_2-ar_1,\ r_3-r_1}
\begin{pmatrix}
1 & -1 & 1 & -1 \\
0 & a+1 & -a-1 & a+1 \\
0 & a+1 & 0 & 2
\end{pmatrix}
\xrightarrow{r_3-r_2}
\begin{pmatrix}
1 & -1 & 1 & -1 \\
0 & a+1 & -a-1 & a+1 \\
0 & 0 & a+1 & 1-a
\end{pmatrix}.
\]

讨论：
\begin{itemize}
    \item 当 \(a \neq -1\) 时：\(\operatorname{rank}(A) = 3 = \operatorname{rank}(A \mid \boldsymbol{b})\)，有唯一解：
    \[
    \begin{aligned}
    x_3&=\frac{1-a}{a+1},\\
    x_2&=1+x_3=\frac{2}{a+1},\\
    x_1&=-1+x_2-x_3=0.
    \end{aligned}
    \]
    \item 当 \(a = -1\) 时：
    \[
    (A \mid \boldsymbol{b}) \to \begin{pmatrix}
    1 & -1 & 1 & -1 \\
    0 & 0 & 0 & 0 \\
    0 & 0 & 0 & 2
    \end{pmatrix},
    \]
    \(\operatorname{rank}(A) = 1\)，\(\operatorname{rank}(A \mid \boldsymbol{b}) = 2\)，无解。
\end{itemize}

(ii) 计算特征多项式。由于 \(A\) 是实对称矩阵，总是可对角化的。
\[
|A - \lambda I| = \begin{vmatrix}
1-\lambda & a & 1 \\
a & 1-\lambda & -1 \\
1 & -1 & 1-\lambda
\end{vmatrix}
= -(\lambda-1-a)\bigl[\lambda^2+(a-2)\lambda-(a+1)\bigr].
\]

特征值为
\[
\lambda_1 = 1+a,\quad
\lambda_{2,3} = \frac{2-a \pm \sqrt{a^2+8}}{2}.
\]

对一切 \(a \in \mathbb{R}\)，三个特征值均为实数；其中
\(\lambda_2\ne\lambda_3\)，而 \(\lambda_1\) 与另一个特征值重合当且仅当
\(a=\pm1\)。无论特征值是否有重根，实对称矩阵均正交相似于对角阵，故 \(A\) 对一切实数 \(a\) 均可对角化。

\boxed{\text{注}} 当 \(a=-1\) 时，特征值为 \(0,0,3\)，故
\(\operatorname{rank}(A)=1\)，与方程组讨论中的秩骤降一致；但增广矩阵秩为 \(2\)，所以非齐次方程组无解。

\fi%
\item \textbf{（矩阵多项式+对角化）} 设 \(A = \begin{pmatrix} 1 & -3 & 3 \\ 3 & -5 & 3 \\ 6 & -6 & 4 \end{pmatrix}\)。

\begin{enumerate}[label=(\roman*)]
    \item 求 \(A\) 的特征值与特征向量，判断 \(A\) 是否可对角化；
    \item 计算 \(A^{100}\)。
\end{enumerate}

\ifshowsolutions%
\boxed{\textbf{解}}
(i) 计算特征多项式：
\begin{align*}
|A-\lambda I| &= \begin{vmatrix}
1-\lambda & -3 & 3 \\
3 & -5-\lambda & 3 \\
6 & -6 & 4-\lambda
\end{vmatrix}
\xrightarrow{r_3-2r_2}
\begin{vmatrix}
1-\lambda & -3 & 3 \\
3 & -5-\lambda & 3 \\
0 & 4+2\lambda & -2-\lambda
\end{vmatrix} \\[4pt]
&= (2+\lambda)\begin{vmatrix}
1-\lambda & -3 & 3 \\
3 & -5-\lambda & 3 \\
0 & 2 & -1
\end{vmatrix}.
\end{align*}

按第三行展开：
\[
|A-\lambda I| = (2+\lambda)\left[-2\begin{vmatrix} 1-\lambda & 3 \\ 3 & 3 \end{vmatrix} + (-1)\begin{vmatrix} 1-\lambda & -3 \\ 3 & -5-\lambda \end{vmatrix}\right]
\]
\[
= (2+\lambda)\big[-2(3(1-\lambda)-9) - ((1-\lambda)(-5-\lambda) - (-9))\big]
\]
\[
= (2+\lambda)\big[-2(-3\lambda-6) - (\lambda^2+4\lambda+4)\big]
= (2+\lambda)(6\lambda+12 - \lambda^2 - 4\lambda - 4)
\]
\[
= (2+\lambda)(-\lambda^2 + 2\lambda + 8) = -(\lambda+2)^2(\lambda-4).
\]

故特征值为 \(\lambda_1 = 4\)（单根），\(\lambda_2 = \lambda_3 = -2\)（二重根）。

求 \(\lambda_1 = 4\) 的特征向量：解 \((A-4I)\boldsymbol{x} = \boldsymbol{0}\)，
\[
A-4I = \begin{pmatrix} -3 & -3 & 3 \\ 3 & -9 & 3 \\ 6 & -6 & 0 \end{pmatrix}
\to \begin{pmatrix} 1 & 1 & -1 \\ 0 & -12 & 6 \\ 0 & -12 & 6 \end{pmatrix}
\to \begin{pmatrix} 1 & 1 & -1 \\ 0 & 2 & -1 \\ 0 & 0 & 0 \end{pmatrix},
\]
得 \(\boldsymbol{p}_1 = (1, 1, 2)^{\mathsf{T}}\)。

求 \(\lambda = -2\) 的特征向量：解 \((A+2I)\boldsymbol{x} = \boldsymbol{0}\)，
\[
A+2I = \begin{pmatrix} 3 & -3 & 3 \\ 3 & -3 & 3 \\ 6 & -6 & 6 \end{pmatrix}
\to \begin{pmatrix} 1 & -1 & 1 \\ 0 & 0 & 0 \\ 0 & 0 & 0 \end{pmatrix},
\]
基础解系：\(\boldsymbol{p}_2 = (1, 1, 0)^{\mathsf{T}}\)，\(\boldsymbol{p}_3 = (-1, 0, 1)^{\mathsf{T}}\)。

三个特征向量线性无关，故 \(A\) 可对角化。

(ii) 令 \(P = (\boldsymbol{p}_1, \boldsymbol{p}_2, \boldsymbol{p}_3)\)，则 \(A = P \cdot \operatorname{diag}(4, -2, -2) \cdot P^{-1}\)。

求 \(P^{-1}\)：
\[
\det P = \begin{vmatrix} 1 & 1 & -1 \\ 1 & 1 & 0 \\ 2 & 0 & 1 \end{vmatrix} = 2.
\]
伴随矩阵法求逆：
\[
P^{-1} = \frac{1}{2}\begin{pmatrix}
1 & -1 & 1 \\
-1 & 3 & -1 \\
-2 & 2 & 0
\end{pmatrix}.
\]

令 \(\alpha = 4^{100}\)，\(\beta = (-2)^{100} = 2^{100}\)，则
\[
\begin{aligned}
A^{100}
&=P\operatorname{diag}(\alpha,\beta,\beta)P^{-1}\\
&= \frac{1}{2}\begin{pmatrix}
\alpha+\beta & -\alpha+\beta & \alpha-\beta \\
\alpha-\beta & -\alpha+3\beta & \alpha-\beta \\
2(\alpha-\beta) & 2(-\alpha+\beta) & 2\alpha
\end{pmatrix}.
\end{aligned}
\]
其中 \(\alpha = 4^{100} = 2^{200}\)，\(\beta = 2^{100}\)。

\boxed{\text{注}} 此题综合了特征值计算、对角化判定和高次幂计算，是考研的经典综合题型。

\fi%
\item \textbf{（二次型标准形+正交变换+特征值）} 已知二次型
\[
f(x_1, x_2, x_3) = x_1^2 + x_2^2 + x_3^2 - 4x_1x_2 - 4x_1x_3 + 2x_2x_3.
\]

\begin{enumerate}[label=(\roman*)]
    \item 用配方法将 \(f\) 化为标准形，并写出所作的可逆线性变换；
    \item 指出 \(f = 1\) 表示何种二次曲面。
\end{enumerate}

\ifshowsolutions%
\boxed{\textbf{解}}
(i) 配方法：
\begin{align*}
f &= x_1^2 - 4x_1x_2 - 4x_1x_3 + x_2^2 + 2x_2x_3 + x_3^2 \\
  &= (x_1 - 2x_2 - 2x_3)^2 - (2x_2+2x_3)^2 + x_2^2 + 2x_2x_3 + x_3^2 \\
  &= (x_1 - 2x_2 - 2x_3)^2 - 4x_2^2 - 8x_2x_3 - 4x_3^2 + x_2^2 + 2x_2x_3 + x_3^2 \\
  &= (x_1 - 2x_2 - 2x_3)^2 - 3x_2^2 - 6x_2x_3 - 3x_3^2 \\
  &= (x_1 - 2x_2 - 2x_3)^2 - 3(x_2 + x_3)^2.
\end{align*}

令
\[
\begin{cases}
y_1 = x_1 - 2x_2 - 2x_3, \\
y_2 = x_2 + x_3, \\
y_3 = x_3,
\end{cases}
\quad \text{即} \quad
\begin{pmatrix} y_1 \\ y_2 \\ y_3 \end{pmatrix} =
\begin{pmatrix}
1 & -2 & -2 \\
0 & 1 & 1 \\
0 & 0 & 1
\end{pmatrix}
\begin{pmatrix} x_1 \\ x_2 \\ x_3 \end{pmatrix}.
\]

变换矩阵可逆，在此变换下：
\[
f = y_1^2 - 3y_2^2.
\]

(ii) \(f = 1\) 即 \(y_1^2 - 3y_2^2 = 1\)，这是 \(y_1y_2\) 平面上的双曲线沿 \(y_3\) 轴平移而成的双曲柱面。由于二次型的秩为 \(2\)（一个零特征值），正惯性指数为 \(1\)，负惯性指数为 \(1\)，故该曲面为双曲柱面。

\boxed{\text{注}} 配方法是化二次型为标准形的基本方法。虽然本题也可用正交变换法，但配方法在仅需标准形而无需正交变换时更简捷。

\fi%
\item \textbf{（\(A^2 = A\) 型证明+对角化+秩）} 设 \(A\) 为 \(n\) 阶方阵且 \(A^2 = A\)（幂等矩阵）。证明：

\begin{enumerate}[label=(\roman*)]
    \item \(A\) 的特征值只能为 \(0\) 或 \(1\)；
    \item \(A\) 可对角化，且 \(\operatorname{rank}(A) = \operatorname{tr}(A)\)；
    \item \(I - A\) 与 \(A\) 有相同的性质，且 \(\operatorname{rank}(A) + \operatorname{rank}(I - A) = n\)。
\end{enumerate}

\ifshowsolutions%
\boxed{\textbf{解}}
(i) 设 \(\lambda\) 为 \(A\) 的特征值，\(\boldsymbol{x} \neq \boldsymbol{0}\) 为对应特征向量，则
\[
A\boldsymbol{x} = \lambda \boldsymbol{x} \;\Longrightarrow\; A^2\boldsymbol{x} = A(\lambda \boldsymbol{x}) = \lambda^2 \boldsymbol{x}.
\]
由 \(A^2 = A\) 得 \(\lambda^2 \boldsymbol{x} = \lambda \boldsymbol{x}\)，即 \((\lambda^2 - \lambda)\boldsymbol{x} = \boldsymbol{0}\)。因 \(\boldsymbol{x} \neq \boldsymbol{0}\)，\(\lambda^2 - \lambda = 0\)，故 \(\lambda \in \{0, 1\}\)。

(ii) 由 \(A^2 = A\) 知 \(A(I - A) = O\)，故 \(\operatorname{col}(I - A) \subseteq \ker A\)。又对任意 \(\boldsymbol{x} \in \mathbb{R}^n\)：
\[
\boldsymbol{x} = A\boldsymbol{x} + (I - A)\boldsymbol{x},
\]
其中 \(A\boldsymbol{x} \in \operatorname{col}(A)\)，\((I - A)\boldsymbol{x} \in \operatorname{col}(I - A) \subseteq \ker A\)。因此
\[
\mathbb{R}^n = \operatorname{col}(A) + \ker A.
\]

若 \(\boldsymbol{y} \in \operatorname{col}(A) \cap \ker A\)，则 \(\boldsymbol{y} = A\boldsymbol{z}\) 且 \(A\boldsymbol{y} = \boldsymbol{0}\)。而 \(A\boldsymbol{y} = A^2\boldsymbol{z} = A\boldsymbol{z} = \boldsymbol{y}\)，故 \(\boldsymbol{y} = \boldsymbol{0}\)。因此
\[
\mathbb{R}^n = \operatorname{col}(A) \oplus \ker A.
\]

\(\operatorname{col}(A)\) 中每个非零向量均为特征值 \(1\) 的特征向量，\(\ker A\) 中每个非零向量均为特征值 \(0\) 的特征向量。取两组基合并即得 \(n\) 个线性无关的特征向量，故 \(A\) 可对角化。

从而存在可逆矩阵 \(P\) 使
\[
A = P\begin{pmatrix} I_r & 0 \\ 0 & 0 \end{pmatrix} P^{-1},
\]
其中 \(r = \operatorname{rank}(A)\)。特征值 \(1\) 的代数重数（=几何重数）为 \(r\)，故
\[
\operatorname{rank}(A) = r = \text{（特征值 \(1\) 的个数）} = \operatorname{tr}(A).
\]

(iii) 由 \((I - A)^2 = I - 2A + A^2 = I - 2A + A = I - A\)，故 \(I - A\) 也是幂等矩阵，特征值也只能为 \(0\) 或 \(1\)。

由直和分解：\(\operatorname{col}(I - A) = \ker A\)，且 \(\operatorname{col}(A) \cap \operatorname{col}(I - A) = \{\boldsymbol{0}\}\)，故
\[
\operatorname{rank}(A) + \operatorname{rank}(I - A) = \dim(\operatorname{col}(A)) + \dim(\operatorname{col}(I - A)) = n.
\]

\boxed{\text{注}} 幂等矩阵描述沿某个补空间到其像空间的投影，其秩等于迹，且与 \(I-A\) 互补。这组结论适合作为矩阵方程、秩与特征值关系的综合提高训练。

\fi%
\item \textbf{（拓展：Cayley--Hamilton + \(A^n\) + 特征值）} 设 \(A = \begin{pmatrix} 2 & 1 \\ -1 & 0 \end{pmatrix}\)。

\begin{enumerate}[label=(\roman*)]
    \item 用 Cayley-Hamilton 定理求 \(A^n\)（\(n \geq 1\)）；
    \item 利用 \(A^n\) 的表达式，验证 \(A^n\) 的特征值等于 \(A\) 的特征值的 \(n\) 次幂。
\end{enumerate}

\ifshowsolutions%
\boxed{\textbf{解}}
(i) 特征多项式：
\[
|\lambda I - A| = \begin{vmatrix}
\lambda - 2 & -1 \\ 1 & \lambda
\end{vmatrix}
= \lambda(\lambda - 2) + 1 = \lambda^2 - 2\lambda + 1 = (\lambda - 1)^2.
\]

由 Cayley-Hamilton 定理：\((A - I)^2 = O\)，即 \(A^2 = 2A - I\)。

由此递推：
\begin{align*}
A^3 &= A \cdot A^2 = A(2A - I) = 2A^2 - A = 2(2A - I) - A = 3A - 2I, \\
A^4 &= A \cdot A^3 = A(3A - 2I) = 3A^2 - 2A = 3(2A - I) - 2A = 4A - 3I.
\end{align*}

由数学归纳法可得通式：
\[
A^n = nA - (n-1)I,\quad n \geq 1.
\]

验证 \(n=1\)：\(1 \cdot A - 0 \cdot I = A\)，成立。
设 \(n=k\) 成立，则
\[
\begin{aligned}
A^{k+1}
&=A\cdot A^k=A\bigl(kA-(k-1)I\bigr)\\
&=kA^2-(k-1)A
=k(2A-I)-(k-1)A\\
&=(k+1)A-kI.
\end{aligned}
\]
归纳成立。

代入 \(A = \begin{pmatrix} 2 & 1 \\ -1 & 0 \end{pmatrix}\)：
\[
A^n = \begin{pmatrix} 2n & n \\ -n & 0 \end{pmatrix} - \begin{pmatrix} n-1 & 0 \\ 0 & n-1 \end{pmatrix}
= \begin{pmatrix} n+1 & n \\ -n & -n+1 \end{pmatrix}.
\]

(ii) \(A\) 的特征值为 \(\lambda = 1\)（二重根）。\(A^n\) 的特征多项式：
\begin{align*}
|\mu I - A^n| &= \begin{vmatrix}
\mu - (n+1) & -n \\ n & \mu + n - 1
\end{vmatrix} \\
&= (\mu - n - 1)(\mu + n - 1) + n^2 \\
&= \mu^2 - 2\mu - (n^2-1) + n^2 = \mu^2 - 2\mu + 1 = (\mu - 1)^2.
\end{align*}

故 \(A^n\) 的特征值仍为 \(1 = 1^n\)，符合「特征值的 \(n\) 次幂」原理。

\boxed{\text{注}} Cayley-Hamilton 定理将矩阵高次幂转化为多项式低次表示，是计算 \(A^n\) 的又一重要工具。对于特征值有重根的情况，C-H 定理尤其方便。

\fi%
\item \textbf{（拓展：正定矩阵同时对角化）} 设 \(A, B\) 均为 \(n\) 阶实对称矩阵，\(A\) 正定。证明存在可逆矩阵 \(P\) 使得
\[
P^{\mathsf{T}} A P = I,\quad P^{\mathsf{T}} B P = \operatorname{diag}(\mu_1, \mu_2, \ldots, \mu_n),
\]
其中 \(\mu_1, \ldots, \mu_n\) 为 \(A^{-1}B\) 的特征值。进一步，若 \(B\) 也正定，问 \(\mu_i\) 满足什么条件？

\ifshowsolutions%
\boxed{\textbf{解}}
由于 \(A\) 正定，存在可逆矩阵 \(C\) 使得 \(C^{\mathsf{T}} A C = I\)（取 \(C = A^{-1/2}\)）。令 \(B_1 = C^{\mathsf{T}} B C\)，则 \(B_1\) 仍为实对称矩阵。

由实对称矩阵的正交对角化，存在正交矩阵 \(Q\) 使得
\[
Q^{\mathsf{T}} B_1 Q = \operatorname{diag}(\mu_1, \mu_2, \ldots, \mu_n).
\]

令 \(P = C Q\)，则
\[
P^{\mathsf{T}} A P = Q^{\mathsf{T}} (C^{\mathsf{T}} A C) Q = Q^{\mathsf{T}} I Q = I,
\]
\[
P^{\mathsf{T}} B P = Q^{\mathsf{T}} (C^{\mathsf{T}} B C) Q = Q^{\mathsf{T}} B_1 Q = \operatorname{diag}(\mu_1, \ldots, \mu_n).
\]

下证 \(\mu_i\) 为 \(A^{-1}B\) 的特征值。由合同关系 \(P^{\mathsf{T}}BP = \Lambda\) 且 \(P^{\mathsf{T}}AP = I\)，设 \(\boldsymbol{e}_i\) 为标准单位向量，则
\[
B(P\boldsymbol{e}_i) = \mu_i A(P\boldsymbol{e}_i).
\]
两边左乘 \(A^{-1}\)：\(A^{-1}B(P\boldsymbol{e}_i) = \mu_i (P\boldsymbol{e}_i)\)，故 \(\mu_i\) 是 \(A^{-1}B\) 的特征值。

若 \(B\) 也正定，则 \(B_1 = C^{\mathsf{T}} B C\) 合同于 \(B\)（因 \(C\) 可逆），故 \(B_1\) 也正定。正定矩阵的特征值全大于零，故 \(\mu_i > 0\) 对所有 \(i = 1, \ldots, n\)。

\boxed{\text{注}} 同时对角化（又称广义特征值问题）是多元统计分析中的核心理论工具。\(A^{-1}B\) 的实特征值与正定性是此类问题的关键结论。

\fi%
\end{enumerate}

\subsection{证明题专项}

\begin{enumerate}

\item \textbf{（拓展：Sylvester 不等式应用）} 设 \(A\) 为 \(m \times n\) 矩阵，\(B\) 为 \(n \times p\) 矩阵。证明：

\begin{enumerate}[label=(\roman*)]
    \item \(\operatorname{rank}(AB) \geq \operatorname{rank}(A) + \operatorname{rank}(B) - n\)；
    \item 利用 (i) 证明：若 \(A\) 列满秩，则 \(\operatorname{rank}(AB) = \operatorname{rank}(B)\)。
\end{enumerate}

\ifshowsolutions%
\boxed{\textbf{解}}
(i) 方法一（几何法）：将矩阵视为线性映射。\(B: \mathbb{R}^p \to \mathbb{R}^n\)，\(A: \mathbb{R}^n \to \mathbb{R}^m\)。考察限制映射 \(A|_{\operatorname{col}(B)}: \operatorname{col}(B) \to \mathbb{R}^m\)。

由秩—零化度定理：
\[
\dim(\operatorname{col}(B)) = \dim(A(\operatorname{col}(B))) + \dim(\ker A \cap \operatorname{col}(B)).
\]
即 \(\operatorname{rank}(B) = \operatorname{rank}(AB) + \dim(\ker A \cap \operatorname{col}(B))\)。

而 \(\dim(\ker A \cap \operatorname{col}(B)) \leq \dim(\ker A) = n - \operatorname{rank}(A)\)，故
\[
\operatorname{rank}(B) \leq \operatorname{rank}(AB) + n - \operatorname{rank}(A),
\]
即 \(\operatorname{rank}(AB) \geq \operatorname{rank}(A) + \operatorname{rank}(B) - n\)。

(ii) 若 \(A\) 列满秩，则 \(\operatorname{rank}(A) = n\)。由 (i)：
\[
\operatorname{rank}(AB) \geq n + \operatorname{rank}(B) - n = \operatorname{rank}(B).
\]
又因为 \(\operatorname{rank}(AB) \leq \min\{\operatorname{rank}(A), \operatorname{rank}(B)\} = \operatorname{rank}(B)\)，故 \(\operatorname{rank}(AB) = \operatorname{rank}(B)\)。

\fi%
\item \textbf{（\(\operatorname{rank}(A) = \operatorname{rank}(A^{\mathsf{T}}A)\) 证明）} 设 \(A\) 为 \(m \times n\) 实矩阵。证明：

\begin{enumerate}[label=(\roman*)]
    \item \(\operatorname{rank}(A^{\mathsf{T}}A) = \operatorname{rank}(A)\)；
    \item \(\ker(A^{\mathsf{T}}A) = \ker(A)\)；
    \item 对任意 \(\boldsymbol{x} \in \mathbb{R}^n\)，\(A\boldsymbol{x} = \boldsymbol{0} \iff A^{\mathsf{T}}A\boldsymbol{x} = \boldsymbol{0}\)。
\end{enumerate}

\ifshowsolutions%
\boxed{\textbf{解}}
(i) 先证 \(\ker(A^{\mathsf{T}}A) \subseteq \ker(A)\)。设 \(\boldsymbol{x} \in \ker(A^{\mathsf{T}}A)\)，即 \(A^{\mathsf{T}}A\boldsymbol{x} = \boldsymbol{0}\)。两边左乘 \(\boldsymbol{x}^{\mathsf{T}}\)：
\[
\boldsymbol{x}^{\mathsf{T}} A^{\mathsf{T}} A \boldsymbol{x} = 0 \;\Longrightarrow\; \|A\boldsymbol{x}\|^2 = 0 \;\Longrightarrow\; A\boldsymbol{x} = \boldsymbol{0}.
\]
故 \(\boldsymbol{x} \in \ker(A)\)。

反包含 \(\ker(A) \subseteq \ker(A^{\mathsf{T}}A)\) 显然（若 \(A\boldsymbol{x} = \boldsymbol{0}\)，则 \(A^{\mathsf{T}}A\boldsymbol{x} = A^{\mathsf{T}}\boldsymbol{0} = \boldsymbol{0}\)）。

因此 \(\ker(A^{\mathsf{T}}A) = \ker(A)\)。由秩—零化度定理（两矩阵均有 \(n\) 列）：
\[
\operatorname{rank}(A^{\mathsf{T}}A) = n - \dim(\ker(A^{\mathsf{T}}A)) = n - \dim(\ker(A)) = \operatorname{rank}(A).
\]

(ii) 已在上一步中证明，两者零空间相等。

(iii) 由 (ii) 中零空间相等直接得到。注意实矩阵情形下，\(A\boldsymbol{x} = \boldsymbol{0} \iff \|A\boldsymbol{x}\|^2 = \boldsymbol{x}^{\mathsf{T}}A^{\mathsf{T}}A\boldsymbol{x} = 0\) 是核心技巧，利用了实向量范数的正定性。

\boxed{\text{注}} \(A^{\mathsf{T}}A\) 与 \(A\) 同秩是线性代数中最重要的秩等式之一，广泛应用于最小二乘法、广义逆和奇异值分解中。

\fi%
\item \textbf{（相似性与对角化证明）} 设 \(A, B\) 均为 \(n\) 阶方阵，且 \(A \sim B\)（相似）。证明：

\begin{enumerate}[label=(\roman*)]
    \item 若 \(A\) 可逆，则 \(B\) 也可逆，且 \(A^{-1} \sim B^{-1}\)；
    \item 若 \(A\) 可对角化，则 \(B\) 也可对角化；
    \item \(A^{\mathsf{T}} \sim B^{\mathsf{T}}\)。
\end{enumerate}

\ifshowsolutions%
\boxed{\textbf{解}}
由 \(A \sim B\)，存在可逆矩阵 \(P\) 使得 \(B = P^{-1}AP\)。

(i) 若 \(A\) 可逆，则 \(|B| = |P^{-1}AP| = |P^{-1}||A||P| = |A| \neq 0\)，故 \(B\) 可逆。

又 \(B^{-1} = (P^{-1}AP)^{-1} = P^{-1}A^{-1}P\)，故 \(A^{-1} \sim B^{-1}\)。

(ii) 若 \(A\) 可对角化，存在可逆矩阵 \(Q\) 使得 \(Q^{-1}AQ = \Lambda\)（对角阵）。则
\[
\Lambda = Q^{-1}AQ = Q^{-1}(PBP^{-1})Q = (Q^{-1}P)B(P^{-1}Q) = (P^{-1}Q)^{-1}B(P^{-1}Q),
\]
故 \(B\) 也可对角化（相似于同一个对角阵 \(\Lambda\)）。

(iii) \(B^{\mathsf{T}} = (P^{-1}AP)^{\mathsf{T}} = P^{\mathsf{T}}A^{\mathsf{T}}(P^{-1})^{\mathsf{T}} = P^{\mathsf{T}}A^{\mathsf{T}}(P^{\mathsf{T}})^{-1}\)。令 \(R = (P^{\mathsf{T}})^{-1}\)，则
\[
B^{\mathsf{T}} = R^{-1}A^{\mathsf{T}}R,
\]
故 \(A^{\mathsf{T}} \sim B^{\mathsf{T}}\)。

\boxed{\text{注}} 相似关系保持矩阵的大部分重要性质：特征值、行列式、迹、秩、可逆性、可对角化性等。但相似不保持对称性（除非相似变换矩阵是正交矩阵）。

\fi%
\item \textbf{（正定性证明）} 设 \(A\) 为 \(n\) 阶正定矩阵，\(S\) 为 \(n \times m\) 列满秩矩阵（\(m \leq n\)）。证明 \(S^{\mathsf{T}} A S\) 为正定矩阵。并给出反例说明若 \(S\) 不列满秩，结论不成立。

\ifshowsolutions%
\boxed{\textbf{解}}
对称性：\((S^{\mathsf{T}} A S)^{\mathsf{T}} = S^{\mathsf{T}} A^{\mathsf{T}} S = S^{\mathsf{T}} A S\)（\(A\) 对称），故 \(S^{\mathsf{T}} A S\) 为对称矩阵。

正定性：任取非零向量 \(\boldsymbol{x} \in \mathbb{R}^m\)，\(\boldsymbol{x} \neq \boldsymbol{0}\)。由于 \(S\) 列满秩，\(S\boldsymbol{x} \neq \boldsymbol{0}\)。又 \(A\) 正定，故
\[
\boldsymbol{x}^{\mathsf{T}} (S^{\mathsf{T}} A S) \boldsymbol{x}
= (S\boldsymbol{x})^{\mathsf{T}} A (S\boldsymbol{x}) > 0.
\]
因此 \(S^{\mathsf{T}} A S\) 正定。

反例：取 \(A = \begin{pmatrix} 1 & 0 \\ 0 & 1 \end{pmatrix}\)（正定），\(S = \begin{pmatrix} 1 & 2 \\ 0 & 0 \end{pmatrix}\)（不列满秩）。则
\[
S^{\mathsf{T}}AS = \begin{pmatrix} 1 & 0 \\ 2 & 0 \end{pmatrix} \begin{pmatrix} 1 & 2 \\ 0 & 0 \end{pmatrix} = \begin{pmatrix} 1 & 2 \\ 2 & 4 \end{pmatrix}.
\]
对 \(\boldsymbol{x} = (2, -1)^{\mathsf{T}} \neq \boldsymbol{0}\)：
\[
\boldsymbol{x}^{\mathsf{T}}(S^{\mathsf{T}}AS)\boldsymbol{x} = \begin{pmatrix} 2 & -1 \end{pmatrix}\begin{pmatrix} 1 & 2 \\ 2 & 4 \end{pmatrix}\begin{pmatrix} 2 \\ -1 \end{pmatrix} = 4 - 4 = 0,
\]
故 \(S^{\mathsf{T}}AS\) 仅为半正定，不正定。

\boxed{\text{注}} 此题揭示了「二次型经线性变换后保持正定性」的充要条件是变换矩阵列满秩，这一结论在约束优化和多元统计中极为重要。

\fi%
\end{enumerate}
